\chapter*{Wprowadzenie do tomu III}
\addcontentsline{toc}{chapter}{Wprowadzenie do tomu III}
W tomie II poznaliśmy proces Wienera, martyngały w czasie ciągłym,
czasy zatrzymania oraz wariację kwadratową. Teraz zmieniamy punkt
widzenia: zamiast wyłącznie badać procesy, chcemy względem nich
\emph{całkować}. Nasuwa się wzór $\int_0^T H_t\,\dd W_t$, ale
sam zapis nie określa jeszcze znaczenia całki.

Przeszkoda jest fundamentalna. Prawie każda trajektoria procesu Wienera
jest ciągła, ale ma na każdym nietrywialnym przedziale nieskończoną
wariację całkowitą. Klasyczna teoria Riemanna--Stieltjesa nie daje więc
od razu całki dla wszystkich interesujących nas funkcji podcałkowych.
Rozwiązaniem jest wykorzystanie informacji probabilistycznej:
przyrosty procesu Wienera są niezależne, mają średnią zero i znaną
wariancję. Dzięki temu sumy całkowe możemy kontrolować w przestrzeni
$L^2$, nawet gdy nie mają sensownej interpretacji w klasycznym
rachunku całkowym wzdłuż każdej pojedynczej trajektorii.

W rozdziale pierwszym budujemy \emph{całkę Wienera} dla
\emph{deterministycznych} funkcji z $L^2([0,T])$. Jest ona szczególnym
przypadkiem całki Itô, lecz jej prostsza konstrukcja pozwala zobaczyć
istotę izometrii i związku z procesami gaussowskimi. W rozdziale
drugim dopuszczamy losową funkcję podcałkową zależną od przeszłości
procesu. Konstrukcję przeprowadzimy najpierw dla funkcji prostych
przewidywalnych, a następnie przez domknięcie w $L^2$ i lokalizację.

Od początku odróżniamy trzy obiekty: zwykłą całkę względem czasu
$\int_0^T H_t\,\dd t$, całkę stochastyczną
$\int_0^T H_t\,\dd W_t$ oraz całkę \emph{po przestrzeni trajektorii}
$\int F(w)\,\dd\mu_W(w)$. Ostatni zapis jest całką względem miary
Wienera, nie zaś synonimem całki Itô. Związki z całką Feynmana,
wzorem Feynmana--Kaca i obrotem Wicka omówimy w dalszej części tomu.

\paragraph{Założenia stałe.}
Niech $T>0$, $(\Omega,\calF,\Prob)$ będzie przestrzenią probabilistyczną,
$(\calF_t)_{0\le t\le T}$ filtracją spełniającą warunki zwykłe,
a $W=(W_t)_{0\le t\le T}$ rzeczywistym standardowym procesem Wienera
\emph{względem tej filtracji}: $W_0=0$, ma ciągłe trajektorie,
jest adaptowany oraz $W_t-W_s\sim\Normal(0,t-s)$ jest niezależny
od $\calF_s$ dla $0\le s<t\le T$. Ostatni warunek jest mocniejszy
niż samo stwierdzenie, że $W$ jest procesem Wienera jako rodzina
zmiennych losowych; będziemy z niego korzystać przy całkowaniu
losowych funkcji podcałkowych.

Przez $L^p(\Omega)$ rozumiemy $L^p(\Omega,\calF,\Prob)$,
a przez $L^p([0,T])$ przestrzeń względem miary Lebesgue'a.
Wartość całki stochastycznej w danej chwili jest określona z dokładnością
do równości prawie na pewno, a cały proces całkowy będziemy wybierać
w wersji o ciągłych trajektoriach.

